\begin{historical}
The recognition that deterministic systems can exhibit unpredictable behaviour marked a change in scientific thought. Before the twentieth century, classical mechanics—embodied by Newton’s laws—was largely viewed as a complete and exact framework: given initial conditions, future behaviour was presumed computable in principle. 

This view was first challenged by Henri Poincaré’s entry to King Oscar II’s prize competition in 1888. Poincaré submitted a memoir claiming to prove stability of certain three-body configurations. The work was accepted and the prize awarded. But before publication, his colleague Lars Edvard Phragmén found an error in a key argument. The correction reversed the conclusion: the orbits Poincaré had declared stable were in fact wildly unstable. He paid for the recall and reprinting of the original memoir—the cost exceeded the prize money—and the corrected version (1890) became the founding document of what would later be called chaos theory. The very calculation meant to demonstrate cosmic predictability had instead uncovered its impossibility.

In the decades that followed, these ideas lay mostly dormant until the rise of computers in the mid-twentieth century allowed for detailed numerical explorations of nonlinear systems. In 1963, Edward Lorenz demonstrated that a set of three differential equations meant to model atmospheric convection could yield drastically different outcomes from imperceptibly different starting points when he re-ran a simulation with rounded initial conditions. This sensitivity, later termed the \QSOpen{}butterfly effect,\QSClose{} became the signature of what we now call chaos. 

Rather than disorder or randomness, chaos refers to the intrinsic unpredictability found in certain deterministic systems. It highlighted the limitations of prediction when geometry of the solution space allows tiny differences to grow exponentially, defying long-term computation even in conceptually simple scenarios.
\end{historical}
